\subsection{Disguised Difference of Squares} 
Sometimes, we have a difference of squares in which one or both squares is a \makeblank.
\begin{Example}
Fully factor the polynomial $x^4-10x^3+25x^2-36$.
\begin{lpic}{}{12}
\foreach \x/\step in {-1/1. Factor out the GCF,
											-2/2. Factor the polynomial factor,
											-8/3. Factor each new polynomial factor,
											-11/4. Put it all together}
	\regtextleft{0}{\x}{\step};
\end{lpic}
\end{Example}
\noi We can check for this whenever we have 3 terms that are \makeblank.
\begin{Example}
Fully factor the polynomial $x^4-4x^2+12x-9$.
\begin{lpic}{}{11}
\foreach \x/\step in {-1/1. Factor out the GCF,
											-2/2. Factor the polynomial factor,
											-7/3. Factor each new polynomial factor,
											-10/4. Put it all together}
	\regtextleft{0}{\x}{\step};
\end{lpic}
\end{Example}
